%%
% BIThesis 本科毕业设计论文模板 —— 使用 XeLaTeX 编译 The BIThesis Template for Undergraduate Thesis
% This file has no copyright assigned and is placed in the Public Domain.
%%

% 中英文摘要章节
\begin{abstract}
% 中文摘要正文从这里开始
自重组轮足机器人兼具轮式移动效率、腿式姿态调节能力与模块化重构特性，在复杂地形通过、协同运输和多机器人系统扩展等场景中具有重要研究价值。针对多个轮足机器人单体通过铰接机构构成链式系统后存在的非完整运动约束、车体间强耦合以及单体平衡稳定性要求高等问题，本文围绕链式自重组轮足机器人协同控制方法展开研究。

首先，依据单体轮足机器人机构特征建立腿部运动学模型，推导关节空间与等效腿长、腿摆角之间的映射关系，并基于虚拟模型控制方法建立腿部支撑力、摆杆力矩与关节驱动力矩之间的等效关系。其次，针对尾部辅助机构在离地与触地两类支撑工况下对机体动力学的影响，建立单体机器人动力学模型，并在平衡点附近线性化得到状态空间表达；在此基础上设计线性二次型调节器，实现对机体姿态、腿部摆角、前向运动状态和尾部姿态的综合反馈控制。针对腿长变化导致模型参数变化的问题，进一步采用参数化增益拟合方法描述控制器增益与构型参数之间的关系，以提高控制器对变构型过程的适应性。

在链式系统协同控制方面，本文根据相邻单体之间的铰接几何关系建立链式运动学约束模型，并由头车运动指令和各单体航向状态推导后续单体的期望速度传播关系。进一步地，构建分层协同控制框架：上层依据链式运动学模型生成满足几何约束的期望运动序列，下层通过单体平衡控制器实现对期望状态的稳定跟踪。该方法为后续引入分布式模型预测控制约束优化和鲁棒扰动补偿提供了理论基础。研究结果表明，所建立的建模与控制框架能够统一描述单体平衡控制与链式协同运动之间的关系，为链式自重组轮足机器人实机控制和性能优化提供了可行技术途径。

\end{abstract}

% 如需手动控制换行连字符位置，可写 aa\-bb，详见
% https://bithesis.bitnp.net/faq/hyphen.html

% 英文摘要章节
\begin{abstractEn}
% 英文摘要正文从这里开始
Self-reconfigurable wheel-legged robots combine the efficiency of wheeled locomotion, the posture-adjusting capability of legged mechanisms, and the structural flexibility of modular robotic systems. They are therefore promising for rough-terrain traversal, cooperative transportation, and scalable multi-robot applications. This thesis investigates cooperative control of a chain-structured self-reconfigurable wheel-legged robotic system, in which multiple balancing wheel-legged modules are connected through articulated joints. The main challenges arise from nonholonomic motion constraints, strong inter-module coupling, and the stability requirement of each balancing module.

First, the kinematic model of the leg mechanism is established according to the wheel-legged structure, and the mapping between the joint coordinates and the equivalent leg length and leg swing angle is derived. Based on virtual model control, the relationship between the virtual supporting force, hip torque, and joint actuation torque is formulated. Then, considering the influence of the tail-assisted mechanism under different contact conditions, a dynamic model of the single module is developed and linearized around the equilibrium point. A linear quadratic regulator is designed to stabilize the body posture, leg swing motion, longitudinal motion, and tail posture. To account for configuration variation caused by changes in leg length, a parameterized gain fitting method is introduced to describe the dependence of the feedback gain on structural parameters.

For cooperative control of the chain-structured system, the geometric constraints between adjacent modules are analyzed, and the velocity propagation relationship from the leading module to the following modules is derived from the articulated kinematic model. A hierarchical cooperative control framework is further constructed: the upper layer generates feasible reference motion sequences satisfying the chain constraints, while the lower layer tracks the desired states through the balancing controller of each module. The proposed modeling and control framework provides a unified description of single-module stabilization and chain-level cooperative motion, and lays a theoretical foundation for further constrained optimization using distributed model predictive control and robust disturbance compensation.
\end{abstractEn}
